
\textbf{Research Questions:}

\begin{itemize}
\item \textbf{RQ1:} Are SE\_N5 and min\_logprob statistically independent uncertainty signals?
\item \textbf{RQ2:} Does SE contribute predictive signal for correctness beyond min\_logprob?
\item \textbf{RQ3:} Does the cross-model judge design produce low circularity?
\end{itemize}

\subsubsection{4.1 Dataset}

\textbf{TriviaQA dev (rc.nocontext).} The closed-book variant of TriviaQA dev (\texttt{mandarjoshi/trivia\textbackslash \_qa}, \texttt{rc.nocontext}, \texttt{validation} split), first N=2500 prompts (seed=42). TriviaQA is a challenging factual QA benchmark where hallucination is frequent, providing sufficient error diversity for UQ evaluation. The correctness rate as assessed by LM-judge is 34.8\% (870/2500 prompts). The no-context variant isolates parametric knowledge from retrieval artifacts, matching the prior baseline configuration.

The N=2500 evaluation constitutes the primary experimental record. It was preceded by a proof-of-concept run at N=300 that validated pipeline correctness, yielding Pearson |r|=0.049, partial $R^{2}=0.0101$, LRT p=0.1504, Spearman $\rho$(SE, correctness)=$-0.081$, and correctness rate 34.3\% (103/300).

\begin{table}[htbp]
\centering
\resizebox{\columnwidth}{!}{%
\begin{tabular}{ll}
\toprule
Property & Value \\
\midrule
Dataset & TriviaQA dev, rc.nocontext \\
N prompts & 2500, seed=42 \\
Task type & Closed-book short-answer factual QA \\
Answer format & Free-form, $\leq 50$ tokens \\
Correctness rate (LM-judge) & 34.8\% (870/2500) \\
\bottomrule
\end{tabular}
}
\end{table}

\subsubsection{4.2 Models}

\textbf{Generator:} Llama-3.1-8B-Instruct (bfloat16, device\_map=''auto''). N=5 stochastic samples per prompt (temperature=0.7, top\_p=0.95, max\_new\_tokens=50). One greedy decode per prompt for min\_logprob extraction.

\textbf{NLI Backbone:} \texttt{cross-encoder/nli-deberta-v3-small} for bidirectional NLI clustering (batch size 32).

\textbf{LM-Judge:} Qwen-2.5-7B-Instruct (cross-model; Llama outputs evaluated by a different model family; batch size 16).

\subsubsection{4.3 Baselines}

\textbf{min\_logprob (standalone):} The minimum token log-probability from a single greedy decode, achieving AUROC approximately 0.825 on TriviaQA dev (internal pipeline experiment h-m1, not published externally). Reference signal for independence test (RQ1) and conditional LR (RQ2).

\textbf{SE\_N5 (standalone):} Semantic Entropy from N=5 stochastic samples per Kuhn et al. (2023). The alternative standalone signal; focus of independence test against min\_logprob. Standalone AUROC for SE\_N5 is not the focus of this work; this experiment tests independence preconditions, not standalone ranking.

Note: Ensemble AUROC is not evaluated in this work. The experiment tests independence preconditions. Given the null result for partial $R^2$, linear ensemble AUROC improvement is not predicted; follow-up work investigates nonlinear combinations.

\subsubsection{4.4 Evaluation Metrics and Hyperparameters}

\begin{table*}[htbp]
\centering
\resizebox{\dimexpr 2\columnwidth+\columnsep\relax}{!}{%
\begin{tabular}{lllll}
\toprule
Metric & Purpose & Pre-registered Threshold &  &  \\
\midrule
Pearson \ & r\ & (SE\_N5, min\_logprob) & Primary independence metric & < 0.70 (PASS); > 0.85 (ABANDON) \\
Partial $R^{2}(SE)$ in conditional LR & Conditional predictive contribution & $\geq$ 0.02, LRT p < 0.05 &  &  \\
Spearman $\rho$(SE, judge\_correctness) & Circularity diagnostic & \ & $\rho$\ & < 0.40 \\
SE variance & Non-degeneracy check & > 0.01 &  &  \\
Fraction degenerate & Non-degeneracy check & = 0 &  &  \\
\bottomrule
\end{tabular}
}
\end{table*}

\begin{table}[htbp]
\centering
\resizebox{\columnwidth}{!}{%
\begin{tabular}{ll}
\toprule
Hyperparameter & Value \\
\midrule
N stochastic samples & 5 \\
Temperature (stochastic) & 0.7 \\
Top-p & 0.95 \\
Max new tokens & 50 \\
NLI batch size & 32 \\
Judge batch size & 16 \\
Pearson r threshold & 0.70 \\
Partial $R^2$ threshold & 0.02 \\
Circularity threshold ($\rho$) & 0.40 \\
ABANDON threshold & 0.85 \\
\bottomrule
\end{tabular}
}
\end{table}

